\chapter{Arquitectura e implementación del software}
\label{ap:arquitectura}

Este apéndice reúne los detalles de implementación del \textit{framework} WoE--AC ---el diagrama de clases y las rutas de los módulos en el repositorio--- que complementan la descripción metodológica del Capítulo~\ref{ch:area-estudio} sin recargar su narrativa. El repositorio organiza el trabajo en una arquitectura modular que integra la rasterización/clasificación de las capturas, el cálculo de Pesos de Evidencia y la simulación con autómatas celulares; la Figura~\ref{fig:system_architecture} muestra el diagrama de clases UML (generado con PlantUML; véase \texttt{diagramas/architecture.puml}).

\begin{figure}[H]
\centering
\includegraphics[width=\textwidth,height=0.72\textheight,keepaspectratio]{figures/architecture.png}
\caption{Arquitectura del framework WoE-AC para modelado de crecimiento urbano. Diagrama UML mostrando las clases del sistema organizadas en paquetes: Configuration (configuración), Data/IO (manejo de datos históricos), Feature \& Spatial (extracción de características), Weight of Evidence (cálculo de pesos), Cellular Automaton (simulación AC) y Orchestrator (pipeline y calibración de parámetros). El diagrama se genera desde el código fuente usando PlantUML (ver \texttt{diagramas/architecture.puml}).}
\label{fig:system_architecture}
\end{figure}

En la rama de \textbf{imágenes}, el flujo opera sobre matrices RGB ya homogeneizadas. El módulo \texttt{FeatureExtractor} (paquete \path{tesis_ac/pipeline/feature\_extraction}) construye las 23 variables resumidas en la Tabla~\ref{tab:features}. La clase \texttt{FeaturePreprocessor} aplica \texttt{StandardScaler} y PCA (ocho componentes por defecto). El componente \texttt{SatelliteImageProcessor} (nombre heredado en el código; procesa matrices RGB de las capturas, véase \path{tesis_ac/pipeline/clustering.py}) agrupa \texttt{ImageClusterer} e \texttt{ImageClassifier}: K-means para $k \in \{2,3,4\}$ y SVM lineal entrenado con pseudoetiquetas. La varianza acumulada explicada por PCA se consulta con \texttt{get\_variance\_explained} del preprocesador. La orquestación de alto nivel está en \texttt{SatelliteImagePipeline} (\path{tesis_ac/pipeline/main\_pipeline.py}).

En la cadena \textbf{WoE--AC} usada en la validación quinquenal de la tesis intervienen \texttt{WoECalculator}, \texttt{WoECellularAutomaton} y la simulación implementada en los módulos \path{tesis_ac/woe/} y \path{tesis_ac/ca/}. La configuración efectiva ---umbral, peso de vecindad y ponderación por IV--- se encuentra registrada en \path{data/processed/quinquenal_best_config.json}. Los experimentos son reproducibles desde la raíz del repositorio con los scripts \texttt{validate\_quinquenal\_*.py}, que leen los mapas en \path{data/processed/standardized_maps/} y el modelo WoE entrenado \path{data/processed/woe_pooled_1984_2010.pkl}; el entrenamiento agregado del WoE se obtiene con \texttt{train\_woe\_pooled.py}. El repositorio conserva, además, un pipeline experimental más amplio en \path{tesis_ac/run_classify_and_grid.py}, que no forma parte de la narrativa defendida en esta tesis (la cual se centra en la búsqueda en cuadrícula y el ajuste manual del umbral descritos en el Capítulo~\ref{ch:modelo-crecimiento-urbano}).
